\subsection{代数中的范数与迹}

定义2. 设 A 是一个 K-代数，且它是有限维自由 K-模。对每个元素 \(a \in \mathbf{A}\) ，a 相对于 A 和 K 的迹（相应地，范数†，相应地，特征多项式）是 K-模 A 的自同态 \(x \mapsto ax\) 的迹（相应地，行列式，相应地，特征多项式）。

 \(a \in A\) 相对于 A 和 K 的迹、范数和特征多项式分别记为 \(\operatorname{Tr}_{\mathbf{A}/\mathbf{K}}(a)\) ， \(\operatorname{N}_{\mathbf{A}/\mathbf{K}}(a)\) 和 \(\operatorname{Pc}_{\mathbf{A}/\mathbf{K}}(a; \mathbf{X})\) ；当没有混淆风险时，我们在此记号中省略 K 甚至 A。注意， \(a \in A\) 的迹（相应地，范数、特征多项式）就是 a 相对于 A-模 \(A_{s}\) 的迹（相应地，范数，特征多项式）。

假设 A 是乘积 \(A_{1} \times A_{2} \times \cdots \times A_{m}\) ，即有限个 K 上的有限维代数（它们均为自由 K-模）的乘积。利用上述注记以及第 2 号的命题 1，对每个元素

\[
a = (a _ {1}, \dots , a _ {m}) \in \mathrm{A},
\]

\[
\mathrm{Tr} _ {\mathbf {A} / \mathbb {K}} (a) = \sum_ {i = 1} ^ {m} \mathrm{Tr} _ {\mathbf {A} _ {i} / \mathbb {K}} (a _ {i}), \qquad \mathrm{N} _ {\mathbf {A} / \mathbb {K}} (a) = \prod_ {i = 1} ^ {m} \mathrm{N} _ {\mathbf {A} _ {i} / \mathbb {K}} (a _ {i})\tag{18}
\]

\[
\mathrm{Pc} _ {\mathbf {A} / \mathbf {K}} (a; \mathbf {X}) = \prod_ {i = 1} ^ {m} \mathrm{Pc} _ {\mathbf {A} _ {i} / \mathbf {K}} (a _ {i}; \mathbf {X}).
\]

类似地，第 2 号的命题 2 表明，若 A 和 A' 是两个代数，它们都是自由 K-模，在 K 上的维数分别为 n 与 \(n'\) ，则对于 \(a \in A\) ， \(a' \in A'\) 。

\begin{enumerate}
\def\labelenumi{(\arabic{enumi})}
\setcounter{enumi}{18}
\tightlist
\item
\end{enumerate}

\[
\operatorname{Tr} _ {\mathbf {A} \otimes \mathbf {A} ^ {\prime}} (a \otimes a ^ {\prime}) = \operatorname{Tr} _ {\mathbf {A}} (a) \operatorname{Tr} _ {\mathbf {A} ^ {\prime}} (a ^ {\prime})\tag{20}
\]

\[
\mathrm{N} _ {\mathbf {A} \otimes \mathbf {A} ^ {\prime}} (a \otimes a ^ {\prime}) = (\mathrm{N} _ {\mathbf {A}} (a)) ^ {n ^ {\prime}} (\mathrm{N} _ {\mathbf {A} ^ {\prime}} (a ^ {\prime})) ^ {n}.
\]

最后，设 A 为 K 上的有限维代数，且它是自由 K-模，h 为 K 到交换环 \(K'\) 的一个同态， \(A' = A_{(K')}\) 为通过 h 作标量扩张而由 A 得到的 \(K'\) -代数。由第 1 号公式 (12) 可知，对所有 \(a \in A\) ，

\[
\begin{array}{r l} \mathrm{Tr} _ {\mathbf {A} ^ {\prime} / \mathbf {K} ^ {\prime}} (1 \otimes a) & = h (\mathrm{Tr} _ {\mathbf {A} / \mathbf {K}} (a)), \quad \mathrm{N} _ {\mathbf {A} ^ {\prime} / \mathbf {K} ^ {\prime}} (1 \otimes a) = h (\mathrm{N} _ {\mathbf {A} / \mathbf {K}} (a)) \\ \mathrm{Pc} _ {\mathbf {A} ^ {\prime} / \mathbf {K} ^ {\prime}} (1 \otimes a; \mathbf {X}) & = \tilde {h} (\mathrm{Pc} _ {\mathbf {A} / \mathbf {K}} (a; \mathbf {X})) \end{array}\tag{21}
\]

其中 \(\tilde{h}\) 是同态 \(\mathbf{K}[\mathbf{X}] \to \mathbf{K}'[\mathbf{X}]\) ，它由 \(h\) 导出。特别地，对于 \(\mathbf{K}' = \mathbf{K}[\mathbf{X}]\) ，记 \(\mathrm{A}[\mathbf{X}] = \mathrm{A} \otimes_{\mathbf{K}} \mathbf{K}[\mathbf{X}]\) ，我们得到

\[
\mathrm{Pc} _ {\mathbf {A} / \mathbf {K}} (a; \mathbf {X}) = \mathrm{N} _ {\mathbf {A} [ \mathbf {X} ] / \mathbf {K} [ \mathbf {X} ]} (\mathbf {X} - a).\tag{22}
\]

更一般地，若 \(\mathbf{K}'\) 是交换 \(\mathbf{K}\) -代数且 \(\mathbf{A}' = \mathbf{A} \otimes_{\mathbf{K}} \mathbf{K}'\) ，则对所有 \(x \in \mathbf{A}'\) ，

\[
\mathrm{Pc} _ {\mathbf {A} / \mathbf {K}} (a; x) = \mathrm{N} _ {\mathbf {A} ^ {\prime} / \mathbf {K} ^ {\prime}} (x - a).
\]

例子。(1) 设 A 是 K 上类型为 \((\alpha, \beta)\) 的二次代数， \((e_1, e_2)\) 是一个类型为 \((\alpha, \beta)\) 的基（§ 2，第 3 号）。对于 \(x = \xi e_1 + \eta e_2\) ， \(\operatorname{Tr}_{A/K}(x) = 2\xi + \beta \eta\) 和 \(\mathrm{N}_{A/K}(x) = \xi^2 + \beta \xi \eta - \alpha \eta^2\) ；因此这些函数与 \(x\) 的凯莱迹和凯莱范数相同（§ 2，第 24 号）。

\begin{enumerate}
\def\labelenumi{(\arabic{enumi})}
\setcounter{enumi}{1}
\item
  设 A 是 K 上的一个四元数代数。直接计算可以验证 \(\mathrm{Tr}_{\mathbf{A} / \mathbb{K}}(x) = 2\mathrm{T}(x)\) 和 \(\mathrm{N}_{\mathbf{A} / \mathbb{K}}(x) = (\mathrm{N}(x))^2\) ，其中 \(\mathbf{T}\) 和 \(\mathbf{N}\) 是凯莱迹和凯莱范数（§ 2，第 4 号）。
\item
  设 \(\mathbf{A} = \mathbf{M}_n(\mathbf{K})\) ，并令 \((E_{ij})\) （ \(\mathbf{A}\) 的典范基，II，§ 10，第 3 号）按字典序排列。则立即可见，对每个矩阵 \(\mathbf{X} = \sum_{i,j} \xi_{ij} E_{ij}\) ，（ \(n^2\) 阶的）自同态 \(Y \mapsto XY\) 的矩阵具有形式
\end{enumerate}

\[
\left( \begin{array}{c c c c} X & 0 & \ldots & 0 \\ 0 & X & \ldots & 0 \\ . & . & . & . \\ 0 & 0 & \ldots & X \end{array} \right)
\]

，因此 \(\operatorname{Tr}_{\mathbf{A} / \mathbb{K}}(X) = n.\operatorname{Tr}(X)\) 和 \(\mathrm{N}_{\mathbf{A} / \mathbb{K}}(X) = (\det (X))^n\)

